В качестве практической части курсовой работы я самостоятельно реализовала рассмотренные мною в теоретической части работы аллокаторы: линейный, пулловый и стековый.

\subsection{Линейный аллокатор}
\inputminted{cpp}{linear_allocator.h}

Данный аллокатор хранит в себе адрес начала большого непрерывного блока памяти, который ему отдан. При инициализации аллокатора этот блок выделяется с помощью функции \texttt{malloc} языка C, а при уничтожении "--- освобождается с помощью функции \texttt{free}.

Помимо этого, у линейного аллокатора есть поля \texttt{used} и \texttt{end} "--- первое из них отвечает за объём памяти, которая уже была выделена, а второе "--- за суммарный объём памяти, доступной аллокатору изначально. Эти поля используется для вычисления адреса, который отдаётся пользователю при аллокации, а также для проверки, хватает ли в аллокаторе памяти для выделения.

Метод \texttt{deallocate} в линейном аллокаторе не имеет смысла, а \texttt{reset} устанавливает \texttt{used} в $0$, тем самым говоря аллокатору, что вся память была освобождена и может быть переиспользована.

\subsection{Стековый аллокатор}
\inputminted{cpp}{stack_allocator.h}

Стековый аллокатор по своей реализации схож с линейным, он включает в себя те же поля \texttt{start}, \texttt{used} и \texttt{end}, а также имеет аналогичный механизм выделения начального блока памяти. Однако в дополнение к основным полям стековый аллокатор имеет поле \texttt{stack_pointer}, реализованное на односвязном списке. При аллокации в список добавляется новый узел, в который записывается информация о размере выделенной памяти; этот узел становится новой вершиной стека.

Деаллокация возможна только в том случае, если переданный указатель совпадает с предыдущей вершиной стека, то есть он получается из текущей вершины стека смещением влево на размер последнего выделенного блока. В таком случае верхний элемент удаляется из стека, а \texttt{used} уменьшается на размер блока. \texttt{reset} полностью очищает стек, устанавливая указатель на его вершину в \texttt{nullptr} и устанавливая \texttt{used} в $0$.

\subsection{Пулловый аллокатор}
\inputminted{cpp}{pool_allocator.h}

Как и все предыдущие аллокаторы, пулловый аллокатор хранит адрес памяти, изначально отданной аллокатору. Он также содержит специфичные поля \texttt{block_count} и \texttt{block_size}, которые отвечают за количество блоков, на которые поделена память, и за размер каждого из них. Односвязный список \texttt{free_list} содержит адреса свободных блоков. В конструкторе пуллового аллокатора не только происходит выделение памяти, но и выполняется итерация по количеству блоков: на каждом шаге вычисляется адрес очередного блока, на которые разделена память, и он записывается в список свободных блоков.

При аллокации из списка берётся первый узел, его адрес возвращается пользователю, а указатель на начало списка сдвигается. При деаллокации пользовательский адрес проверяется на корректность (смещение должно быть кратно размеру блока), при успешном выполнении проверки данный адрес добавляется в список свободных и может быть переиспользован.
